\section{Decoding}\label{sec:decoding}

The final QEC building block we'll discuss is decoding.
Suppose we have a circuit $C$ and noise model $\FFF$.
We can run this circuit and look at the resulting syndrome -- the list of violated detectors on this particular run.
We saw in \Cref{subsec:detecting_webs_detect_errors} how to tell if an error violates an individual detector.
But now we have a whole set of violated detectors, and we have to find errors that would have violated exactly these detectors and no others.

To make life more complicated (perhaps counterintuitively), we don't even actually have to find the exact error that occurred!
We say an error $F$ on circuit $C$ is \textit{trivial} iff $C^F \propto C$, i.e.\ if they occur, they don't actually harm our computation.
By definition, they also don't violate any detectors.
%We showed two errors $F$ and $F'$ are equivalent iff $F' = FT$ for some trivial error $T$.
So it makes no difference at all (to either the syndrome nor the overall computation) whether error $F$ occurs or an \textit{equivalent error} $F' = FT$, where $T$ is some trivial error.
This property of different equivalent errors having the same syndrome is called \textit{degeneracy}.
Letting $\TTT$ denote the set of all trivial errors on $C$, what we're really looking for is the most likely coset $F\TTT$ that the actual error belonged to.
A \textit{decoder} is any algorithm that aims to tackle this problem.

In general, finding the most likely error coset is a harder problem than finding the most likely error itself, so many decoders in fact do the latter, given that it's a reasonable proxy for the former.
Indeed, both of the decoders we'll look at in this section take this approach.

\subsection{Decoding graph}

Before we can look at examples of decoders, we discuss what exactly is input to a decoder: the \textit{decoding graph}.
Suppose we're given a circuit $C$ and an error model $\FFF$.
The decoding graph for this $C$ and $\FFF$ is a bipartite graph $G = ((V_\DDD, V_\FFF), E)$.
One set of vertices $V_\DDD$ is in bijection with the group $\DDD$ of all detectors in the circuit; we label these vertices by $(D, e)$, where $D$ is the detector (a \textit{formal} product $D = m_1 \ldots m_r$ of measurement outcomes) and $e = \pm 1$ the value it evaluates to in the absence of noise (the \textit{actual} product $m_1 \ldots m_r = \pm 1$ of measurement outcomes, which is deterministic in the absence of noise).
The other set of vertices $V_\FFF$ is in bijection with the set of all possible errors in $\FFF$.
That is, for each probability distribution $\FFF_j = \{(F_{j,k}, p_{j,k})\}_k$ in $\FFF$, we add a vertex labelled $(F_{j,k}, p_{j,k})$.
An edge connects detector vertex $(D, e)$ with error vertex $(F, p)$ iff error $F$ would violate detector $D$.
Vertices corresponding to errors that violate no detectors can be removed from the graph.
For example, we show the decoding graph for the circuit consisting of consecutive $Z_1 Z_2$ measurements under phenomenological noise:
\begin{equation*}
    \tikzfig{qubits/prelims/decoding/meas_zz_twice_phenom_decoding_graph}
\end{equation*}
This circuit has one detector, so the decoding graph has one detector vertex.
This is labelled $(m_a m_b, e)$, where $m_a = (-1)^a$ and $m_b = (-1)^b$, and we know $m_a m_b = 1$ in the absence of any errors, so $e = 1$.
There are eight error locations, denoted using dashed lines in both the circuit and the decoding graph.
In total there are 28 possible errors -- though 8 of these are the identity so can be omitted from the graph since they can never violate any detectors.
A simplified version is therefore:
\begin{equation}
    \tikzfig{qubits/prelims/decoding/meas_zz_twice_phenom_decoding_graph_simpler}
\end{equation}

This is the `abstract' decoding graph, if you like.
Then on any given run of the circuit $C$, we create a copy $G'$ of $G$ that also records the syndrome (set of detectors that were actually violated during this run).
We do this by updating the labels of all detector vertices from $(D, e)$ to $(D, eD)$, where (confusingly, I admit) the first element $D = m_1 \ldots m_r$ is a formal product, while the second element $eD = \pm 1$ the result of multiplying the expected value $e = \pm 1$ with the actual evaluation $D = m_1 \ldots m_r = \pm 1$ obtained on this run of the circuit.
As a result, vertices labelled $(D, -1)$ in $G'$ correspond to violated detectors, while vertices labelled $(D, 1)$ correspond to non-violated detectors.
\begin{equation}
    \tikzfig{qubits/prelims/decoding/meas_zz_twice_phenom_decoding_graph_simpler_violated}
\end{equation}

This $G'$ can then be given as input to a decoder.
In the next subsection, we'll look at one such decoder, called \textit{belief propagation}.
Other decoders can only be applied to decoding graphs satisfying certain extra constraints; we will shortly look at \textit{minimum weight perfect matching}, which falls into this category.

\subsection{Belief propagation}

One widely used classical decoding algorithm is called \textit{belief propagation} (BP).
This takes as input a classical decoding graph $G'$, which is a bipartite graph in which one set of vertices correspond to errors and their probabilities of occurring, and the other to \textit{parity checks} (classical detectors, essentially) and whether they were violated (the syndrome).
Belief propagation works by iteratively passing messages around this graph, using these to update the probabilities of different errors occurring given that this particular syndrome has occurred, until it's reasonably confident that it has identified the most likely set of errors to have caused this syndrome (hence the name!).

Since the classical decoding graph is the same kind of object as the decoding graph we got from our quantum circuit, we can apply belief propagation to this too!
The bad news is, this doesn't work anywhere near as well as in the classical case, owing to the fact that quantum circuits are almost always degenerate, as described above.
Without going into any details around why BP struggles with degeneracy, the net result is that if we run belief propagation on a quantum decoding graph, the probabilities do start to get updated, but they don't generally converge to a point where the decoder is reasonably confident it has found the most likely set of errors to have caused this syndrome.
All is not lost, however.
We can still use belief propagation as a preprocessing algorithm to update our probabilities initially before passing the updated decoding graph to another decoder -- for example, a \textit{minimum weight perfect matching} decoder.

\subsection{Minimum weight perfect matching}

Suppose we have a decoding graph $G = ((V_\DDD, V_\FFF), E)$ in which every error vertex $(F, p) \in V_\FFF$ is incident to at most two vertices.
Let's augment this with a boundary detector vertex labelled $(\boundary, e_{\boundary})$, where $\boundary$ is the formal product $\boundary = \prod_{D \in \DDD} D$ of all detectors in the circuit, and $e_{\boundary}$ is the expected evaluation of this in the absence of noise, which (again, confusingly -- apologies!) uses the same notation: $e_{\boundary} = \prod_{D \in \DDD} D = \pm 1$.
This value for $e_{\boundary}$ is chosen so that there will always be an even number of violated detectors total.
Let's connect this boundary vertex to all error vertices that are only incident to one edge, so that now all error vertices connect to exactly two edges.
This means we can actually just think of them as edges rather than vertices!
Let's do exactly this, and call the resulting graph $G_\boundary = (V_\DDD \cup \{(\boundary, e_\boundary)\}, V_\FFF)$.
We show this process for the circuit that measures $Z_1 Z_2$ under a noise model that only allows measurement errors:
\begin{equation}
    \tikzfig{qubits/prelims/decoding/meas_zz_thrice_decoding_graph_boundary}
\end{equation}

So now the decoding problem is: find the minimum weight set of edges $\{(F_k, p_k)\}_k$ whose \textit{boundary} is the set of vertices $\{(D, s) : s = -1\}$ corresponding to violated detectors.
Here the weight of a set of edges is $wt(\{(F_k, p_k)\}_k) = \sum_k -log(p_k)$.
We hope the meaning of the boundary of a set of edges is intuitively clear -- formally we mean boundary in the \textit{simplicial complex} sense, where we identify sets with \textit{$\mathbb{F}_2$-chains}, but if this means nothing to you then no problem.

Happily, we can at this point leverage standard graph theoretic algorithms to find such a set!
One popular algorithm is minimum weight perfect matching (MWPM), the details of which we won't go into.
Unlike belief propagation, this algorithm works perfectly well in the quantum setting.
\teague{More to say?}

\subsection{Reducing error vertex degree}

The drawback of the minimum weight perfect matching decoder is that it requires a decoding graph whose error vertices have degree at most two.
Most decoding graphs fail to satisfy this.
However, we can efficiently transform many decoding graphs into a new decoding graph with this property, at the cost of losing some information about the probabilities of errors occurring.
The loss of this information ultimately reduces decoding accuracy.
We won't give an exact procedure here, but will give a quick example, and refer the reader wanting full details to TODO.
Consider the case below, where an error vertex of degree 4 is removed and the probabilities $p_1$ and $p_3$ are updated:
\begin{equation}
     \tikzfig{qubits/prelims/decoding/remove_hyperedge}
\end{equation}
The new probability $p_1'$ is the probability that either $f_1$ occurred, or $f_2 f_3$, but not both: $p_1' = p_1 + p_2 p_3 - p_1 p_2 p_3$.
Likewise for $p_3'$.
In the decoding graph on the left, the violations of $D_1$ and $D_2$, and of $D_3$ and $D_4$, are dependent events -- the error $f_2$ violates all four of them.
In the decoding graph on the right, however, the violations of $D_1$ and $D_2$, and of $D_3$ and $D_4$, are treated as independent events.

\subsection{Belief propagation as a pre-processor}

Applying the degree reduction process hinted at above lets us use minimum weight perfect matching, but at the cost of a loss of accuracy.
We can mitigate this effect by first running belief matching on our decoding graph $G'$, then applying the transformation above, and then applying MWPM.
This whole pipeline is referred to as \textit{belief matching}.
Intuitively, this works because belief matching partially recovers information that's lost when we delete error vertices with degree greater than two.
For example, given the decoding graph above, suppose on a given run all four detectors are violated.
Running BP updates the probabilities associated with each error, and on this particular run will increase the probability $p_2 \mapsto p_2'$ significantly.
\begin{equation}
    \tikzfig{qubits/prelims/decoding/four_detector_decoding_graph_all_violated_bp}
\end{equation}
When we now delete the degree four error vertex update the probabilities of $p_1'$ and $p_3'$, we encode some of this dependency between the violations of $D_1$ and $D_2$, and of $D_3$ and $D_{4}$, into the updated probabilities (i.e.\ $p_1'' = p_1' + p_2' p_3' - p_1' p_2' p_3'$, similar for $p_3''$).
\begin{equation}
    \tikzfig{qubits/prelims/decoding/four_detector_decoding_graph_all_violated_bp_reduced}
\end{equation}
